% !TeX encoding = UTF-8
% !TeX program = xelatex
% 编译两次：xelatex main.tex
\documentclass[UTF8,a4paper,zihao=-4,fontset=fandol]{ctexart}
\usepackage{course-report}

% ==================== 只需先修改这里 ====================
\renewcommand{\ReportTitle}{在此填写课程报告题目}
\renewcommand{\ReportType}{课程报告} % 可改为：实验报告、课程论文、课程设计
\renewcommand{\CourseName}{在此填写课程名称}
\renewcommand{\StudentName}{在此填写姓名} % 仅用于 PDF 作者信息
% ======================================================

\begin{document}
\pagenumbering{Roman}
\MakeReportContents % 第一页直接从目录开始。
\StartReportBody

\section{引言}
\subsection{研究背景与意义}
在此介绍课程主题、研究背景和相关工作，并说明本报告的研究意义。

\subsection{研究目标与内容}
在此说明需要解决的问题、预期目标以及报告的主要内容。

\section{理论基础与研究方法}
\subsection{基本原理}
在此介绍相关概念、理论依据和必要的假设。首次使用缩写或符号时应说明含义。

\subsection{方法与技术路线}
在此说明采用的方法、处理步骤和技术路线，使他人能够理解并复现分析过程。
需要数学表达时，可使用带编号的公式，例如：
\begin{equation}
  \bar{x}=\frac{1}{n}\sum_{i=1}^{n}x_i .
  \label{eq:mean}
\end{equation}
式~\eqref{eq:mean} 中，$n$ 为样本数，$x_i$ 为第 $i$ 个观测值。

\section{数据与实验设计}
\subsection{数据来源与处理}
在此说明数据来源、采集时间、规模、单位、预处理方法和质量控制措施。

\subsection{实验环境与参数}
在此记录软件版本、实验条件和主要参数。表~\ref{tab:settings} 提供可直接修改的三线表。
\begin{table}[htbp]
  \centering
  \caption{实验设置}
  \label{tab:settings}
  \begin{tabularx}{0.9\linewidth}{P{3cm}CL}
    \toprule
    项目 & 设置值 & 说明 \\
    \midrule
    数据规模 & 待填写 & 样本数量或空间范围 \\
    运行环境 & 待填写 & 软件版本及硬件条件 \\
    关键参数 & 待填写 & 参数含义与选择依据 \\
    \bottomrule
  \end{tabularx}
\end{table}

\section{结果与分析}
\subsection{主要结果}
在此展示主要结果，结合图表解释现象。图表应有标题、编号和正文引用；
涉及测量值时应注明单位，涉及外部数据时应说明来源。

\subsection{讨论与局限}
在此讨论结果的成因、与预期或已有研究的差异，以及方法和数据的局限。

\section{结论}
在此对应研究目标总结主要发现，说明课程知识的应用与收获，并提出后续改进方向。

% 手工参考文献无需额外的 .bst 或 .bib 文件。
% 在正文相应位置使用 \cite{ref:source}；提交前删除未引用的占位条目。
\clearpage
\phantomsection\addcontentsline{toc}{section}{参考文献}
\begin{thebibliography}{99}
  \bibitem{ref:source} 作者. 文献题名[文献类型]. 出版信息, 年份.
\end{thebibliography}

% 不需要附录时，可删除下列部分。
\clearpage
\appendix
\section{补充材料}
在此放置补充表格、推导过程或必要的源代码。丰富的排版示例见 examples.tex。

\end{document}
